% from ega4-18-fr.tex
\oldpage[IV-4]{112}

\begin{proposition}[18.2.3]
\EGAReaderTarget{ega4.c04.prop.18-2-3}{IV.18.2.3}
Soit $f:X\to Y$ un morphisme fini de préschémas et soit $\sh{B}=f_*(\sh{O}_X)$. Les conditions suivantes sont équivalentes:
\begin{enumerate}
  \item[a)] $f$ est étale.
  \item[a$'$)] $f$ est un morphisme plat de présentation finie, et pour tout $x\in X$, si l'on pose $y=f(x)$, $\sh{O}_{X,x}/\mathfrak m_y\sh{O}_{X,x}$ est un corps, extension finie séparable de $\kappa(y)$.
  \item[b)] $\sh{B}$ est un $\sh{O}_Y$-Module localement libre et pour tout $y\in Y$, $\sh{B}_y\otimes_{\sh{O}_{Y,y}}\kappa(y)$ est une $\kappa(y)$-algèbre finie séparable (donc le composé direct d'un nombre fini de corps, extensions finies séparables de $\kappa(y)$).
  \item[c)] $\sh{B}$ est un $\sh{O}_Y$-Module localement libre, et l'homomorphisme $\operatorname{astr}_{\sh{B}|\sh{O}_Y}:\sh{B}\to\check{\sh{B}}$ (18.2.2) est bijectif.
\end{enumerate}
\end{proposition}

Compte tenu de ce que $f$ est quasi-compact, l'équivalence de $a)$ et $a')$ a déjà été démontrée (17.6.2). Pour prouver le reste de la proposition, on peut se borner au cas où $Y=\Spec(A)$ et $X=\Spec(B)$ sont affines, $B$ étant une $A$-algèbre finie et $\sh{B}=\widetilde B$. Dire que $f$ est un morphisme de présentation finie équivaut alors à dire que $B$ est un $A$-module de présentation finie (1.4.7). Si en outre $f$ est plat, donc $B$ un $A$-module plat, on sait (Bourbaki, \emph{Alg. comm.}, chap. II, \textsection~5, n\textsuperscript{o}~2, cor. 2 du th. 1) que $B$ est un $A$-module projectif, donc $\sh{B}$ un $\sh{O}_Y$-Module localement libre (\emph{loc. cit.}, n\textsuperscript{o}~2, th. 1), et la réciproque est immédiate. D'autre part $f^{-1}(y)$ n'est autre que le spectre de la $\kappa(y)$-algèbre $\sh{B}(y)=\sh{B}_y\otimes_{\sh{O}_{Y,y}}\kappa(y)$, ce qui achève de prouver l'équivalence de $a')$ et $b)$. Pour voir que $b)$ est équivalente à $c)$, notons que la seconde assertion de $b)$ équivaut au fait que l'homomorphisme $\operatorname{astr}_{\sh{B}(y)|\kappa(y)}:\sh{B}(y)\to\check{\sh{B}(y)}$ est bijectif; comme $\sh{B}(y)=\sh{B}_y/\mathfrak m_y\sh{B}_y$ et $\check{\sh{B}(y)}=\check{\sh{B}}_y/\mathfrak m_y\check{\sh{B}}_y$ et que $\sh{B}_y$ et $\check{\sh{B}}_y$ sont des $\sh{O}_{Y,y}$-modules libres, il résulte de (18.2.2.4) et de Bourbaki, \emph{Alg. comm.}, chap. II, \textsection~3, n\textsuperscript{o}~3, cor. de la prop. 6, que l'homomorphisme $\operatorname{astr}_{\sh{B}_y|\sh{O}_{Y,y}}:\sh{B}_y\to\check{\sh{B}}_y$ est lui aussi bijectif; la réciproque étant évidente, cela achève la démonstration.

Lorsqu'une $\sh{O}_X$-Algèbre $\sh{B}$ vérifie les conditions équivalentes $b)$ et $c)$ de (18.2.3), on dit que $\sh{B}$ est une $\sh{O}_X$-Algèbre \emph{finie et étale}. Lorsque $X=\Spec(A)$ est affine et que l'on a donc $\sh{B}=\widetilde B$, où $B$ est une $A$-algèbre, il revient au même, en vertu de (18.2.3), de dire que $\sh{B}$ est une $\sh{O}_X$-Algèbre finie et étale, ou que $B$ est une $A$-algèbre finie et étale (au sens de (17.3.2)).

\begin{corollary}[18.2.4]
\label{IV.18.2.4}
Soit $f:X\to Y$ un morphisme fini et de présentation finie, et posons $\sh{B}=f_*(\sh{O}_X)$. Soit $y$ un point de $Y$. Les conditions suivantes sont équivalentes:
\begin{enumerate}
  \item[a)] Il existe un voisinage ouvert $U$ de $y$ dans $Y$ tel que la restriction $f^{-1}(U)\to U$ de $f$ soit un morphisme étale.
  \item[b)] $\sh{B}_y$ est un $\sh{O}_{Y,y}$-module libre de type fini et $\sh{B}_y\otimes_{\sh{O}_{Y,y}}\kappa(y)$ est une $\kappa(y)$-algèbre séparable.
\end{enumerate}
\end{corollary}

Il est clair que $a)$ entraîne $b)$ en vertu de (18.2.3). D'autre part, $\sh{B}$ est un $\sh{O}_Y$-Module de présentation finie (1.6.3 et 1.4.7), donc, si $\sh{B}_y$ est un $\sh{O}_{Y,y}$-module libre, il existe un voisinage ouvert $U$ de $y$ dans $Y$ tel que $\sh{B}|U$ soit un $\sh{O}_U$-Module localement libre ($0_{\mathrm I}$, 5.2.7); en outre, par hypothèse, l'homomorphisme $\operatorname{astr}_{\sh{B}_y|\sh{O}_{Y,y}}:\sh{B}_y\to\check{\sh{B}}_y$ étant bijectif, il résulte aussi de ($0_{\mathrm I}$, 5.2.7) que l'on peut supposer $U$ choisi de sorte que

